% !PRINTED 665
\weberpage{665}
%<<CONT
wenn $c_1, c_2, \ldots, c_n$ die durch die Kongruenzen

\weq{3}{\omega_1 \equiv c_1,\ \omega_2 \equiv c_2,\ \ldots,\ \omega_n \equiv c_n\ (\mathrm{mod}\ \pi)}

bestimmten Zahlwerte sind. Hierdurch sind die Werte aller ganzen Funktionen von $z$ im Punkt $\mathfrak{P}$ bestimmt. Jede ganze Funktion $\omega$ ist nach dem Modul $z - c$ einem Ausdruck
\[
a_1\omega_1 + a_2\omega_2 + \cdots + a_n\omega_n
\]
mit konstantem Koeffizienten $a$ kongruent und ist dann und nur dann durch $\pi$ teilbar, wenn
\[
a_1c_1 + a_2c_2 + \cdots + a_nc_n = 0
\]
ist. Demnach folgt:

2. \gesperrt{Unter den ganzen Funktionen von $z$ haben nur die durch $\pi$ teilbaren in $\mathfrak{P}$ den Wert 0.}

Da man jede gebrochene Funktion in $\Omega$ so darstellen kann, daß Zähler und Nenner nicht zugleich durch $\pi$ teilbar sind, so ist hierdurch der Wert einer jeden Funktion in $\mathfrak{P}$ bestimmt.

3. \gesperrt{Der Primfaktor $\pi$ heißt durch den Punkt $\mathfrak{P}$ erzeugt. Er kann nicht durch einen von $\mathfrak{P}$ verschiedenen Punkt erzeugt werden.}

Man kann auch umgekehrt aus jedem Primfaktor $\pi$ in $z$ einen Punkt erzeugen.

Ist $z - c$ die durch $\pi$ teilbare Linearfunktion, so gebe man der Funktion $z$ den Wert $c$. Dann nehme man eine durch $\pi$, aber nicht durch $\pi^2$ teilbare Funktion $\varrho$ und setze nach § 185, (5) für jede Funktion $\eta$:

\weq{4}{\eta = a\varrho^m + \eta_1\varrho^{m+1},}

worin $m$ eine ganze rationale Zahl, $\eta_1$ eine Funktion, deren Nenner nicht durch $\pi$ teilbar ist. Der Funktion $\varrho$ erteile man den Wert 0 und der Funktion $\eta$ den Wert
\[
\begin{aligned}
\eta_0 &= \infty, \text{ wenn } m < 0,\\
\eta_0 &= a, \text{ wenn } m = 0,\\
\eta_0 &= 0, \text{ wenn } m > 0.
\end{aligned}
\]

4. \gesperrt{Die so bestimmten Werte aller Funktionen $\eta$ genügen den Bedingungen 1. bis 5. und konstituieren also einen Punkt, der durch das Primfunktional $\pi$ erzeugt heißt.}

Man kann also aus einer einzigen Funktion $z$ und den Primfaktoren der Linearformen $z - c$ alle Punkte ableiten, in denen
%CONT>>
